\documentclass[11pt,a4paper]{article}

% --- Core Packages ---
\usepackage[utf8]{inputenc}
\usepackage[margin=1in]{geometry}
\usepackage{helvet}
\usepackage{microtype}
\usepackage{booktabs}
\usepackage{tabularx}
\usepackage{amsmath}
\usepackage{enumitem}
\usepackage{hyperref}

\hypersetup{
    colorlinks=true,
    linkcolor=black,
    urlcolor=blue
}

\setlist[itemize]{noitemsep, topsep=2pt}
\setlist[enumerate]{noitemsep, topsep=2pt}

% --- Title Information ---
\title{\textbf{\Large TimeLens AI: A Personalized Time-Prediction and Adaptive Scheduling Engine}\\[0.3em]
\large Software Engineering Project Proposal}
\author{
    \textbf{Vansh Jain} (1024030770) \quad 
    \textbf{Devyang} (1024030765) \quad 
    \textbf{Bhavik} (1024030772)\\[0.5em]
    \small Department of Computer Engineering\\
    \small Thapar Institute of Engineering and Technology
}
\date{August 24, 2026}

\begin{document}

\maketitle
\hrule
\vspace{1.5em}

\section{Higher-Order Goal}
This project aims to enhance personal productivity and temporal literacy by replacing conventional, static task planning with an adaptive, data-driven forecasting architecture. Beyond immediate technical execution, it targets a fundamental systemic issue in productivity software: human time misestimation follows personal, measurable, and systemic patterns. Existing task managers treat user-inputted durations as absolute ground truth, completely ignoring historical behavioral signals.

\section{Problem Statement}
Current task management platforms function primarily as passive storage systems, relying on user-provided completion estimates without validation. This static operational model fails due to five critical structural flaws:
\begin{itemize}
    \item \textbf{Unmitigated Estimation Bias:} Users systematically under- or overestimate durations. Planners fail to track or adjust for these individual, category-specific variance patterns.
    \item \textbf{Assumption of Uniform Productivity:} Human cognitive output fluctuates based on circadian rhythms, energy reserves, and continuous workload, whereas planners assume static throughput.
    \item \textbf{Uncaptured Context-Switching Costs:} External interruptions and task switching degrade operational efficiency, yet existing systems fail to factor these delays into future allocations.
    \item \textbf{Unmodeled Deadline Dynamics:} Dynamic responses to approaching deadlines, ranging from output acceleration to behavioral paralysis, are entirely unaccounted for.
    \item \textbf{Static Schedule Degradation:} Schedules remain rigid following initial delays, compounding downstream timeline failures rather than adaptively re-balancing workload.
\end{itemize}

\textit{TimeLens AI} resolves these structural limitations by constructing per-user statistical models that estimate realistic task completion vectors and automatically re-sequence remaining workloads upon data ingestion.

\section{Objectives}
\begin{itemize}
    \item Engine construction for personalized task duration estimation localized per user and task category.
    \item Algorithmic quantification and auto-correction of user-specific estimation bias profiles.
    \item Modeling time-of-day productivity efficiency and cumulative cognitive workload trends.
    \item Integration of quantifiable task interruption and context-switch overhead metrics.
    \item Implementation of a real-time reactive scheduling engine for dynamic workload re-allocation.
\end{itemize}

\section{Methodology \& Solution Approach}

\subsection{System Architecture \& Pipeline}
The end-to-end operational pipeline processes tasks through four core sequential phases:
\begin{enumerate}
    \item \textbf{Ingestion \& Feature Extraction:} Accepts inputs via Client API Gateway. Computes rolling user-category estimation biases, temporal productivity profiles, interruption metrics, and deadline proximity factors.
    \item \textbf{Predictive Inference:} Evaluates extracted vectors using a machine learning engine (transitioning from baseline linear models to gradient-boosted regressors as data density scales) to generate predicted completion durations and confidence intervals.
    \item \textbf{Dynamic Schedule Re-alignment:} Re-sequences and recalibrates open tasks systematically upon trigger events (completion, creation, or delay).
    \item \textbf{Feedback Loop Persistence:} Commits actual completion metrics to an immutable historical repository, continuously refining parameters.
\end{enumerate}

\subsection{Key Model Features}
Features analyzed include: Task Category, Self-Reported Estimate, Historical User-Category Estimation Bias, Start Timestamp (Circadian Vector), Daily Cumulative Workload, Recorded Interruption Volume, and Temporal Distance to Deadline.

\subsection{Personalization \& Cold-Start Strategy}
Models prioritize individual operational history over global heuristics. For zero-history onboarding, the framework applies a population-level prior, continuously blending and shifting weights toward the personalized statistical model as empirical task data accumulates.

\subsection{Technical Infrastructure}
\begin{itemize}
    \item \textbf{Core Services:} RESTful API Architecture with Token-based Authentication.
    \item \textbf{ML Engine:} \texttt{scikit-learn} / \texttt{XGBoost} gradient-boosted regressor models.
    \item \textbf{Data Tier:} Relational Store for tasks and history; Feature Store for rolling statistical profiles.
    \item \textbf{User Interface:} Real-time visualization dashboard tracking historical bias and live schedules.
\end{itemize}

\section{Evaluation Criteria}
Target performance is evaluated strictly through quantitative KPIs:
\begin{itemize}
    \item \textbf{Prediction Accuracy:} Mean Absolute Error (MAE) between predicted vs. actual durations.
    \item \textbf{Bias Reduction:} Variance reduction between user estimates and empirical durations post-correction.
    \item \textbf{Schedule Adherence:} Percentage of tasks executed within the dynamically projected window.
    \item \textbf{Rebuild Latency:} Execution latency of schedule recalibration upon state change.
    \item \textbf{Missed-Deadline Reduction:} Metric shift in breached deadlines against baseline.
\end{itemize}

\textbf{Execution Case Example:} A user estimates $60\text{ mins}$ for a technical report scheduled for the afternoon. Incorporating historical bias and circadian throughput degradation, \textit{TimeLens AI} predicts $95\text{ mins}$, automatically delaying subsequent downstream tasks by $35\text{ mins}$ to maintain schedule integrity.

\section{Scalability \& System Architecture}
\begin{itemize}
    \item \textbf{Decoupled Infrastructure:} Ingestion, inference, and scheduling operate as isolated microservices to ensure horizontal scale.
    \item \textbf{Stateless Execution:} Database optimization through indexed user querying and per-user model training guarantees isolated computational overhead.
\end{itemize}

\section{Resource Allocation}
Implementation relies exclusively on open-source frameworks (\texttt{scikit-learn}, relational DBs) and existing university infrastructure, requiring zero external funding.

\section{Project Scope \& Deliverables}

\subsection{Phase 1 Deliverables}
\begin{itemize}
    \item Task management and authentication API endpoints.
    \item Baseline user-category estimation bias correction engine.
    \item Primary duration prediction endpoint and core telemetry dashboard.
\end{itemize}

\subsection{Phase 2 Deliverables}
\begin{itemize}
    \item Integrated circadian productivity and interruption metrics engine.
    \item Dynamic schedule rebuilding system with deadline-pressure scaling.
    \item Automated persistent learning loop and live analytics reporting.
\end{itemize}

\section{Risk Matrix \& Mitigation Strategies}
\begin{table}[htbp]
\centering
\small
\begin{tabularx}{\textwidth}{>{\bfseries}l X}
\toprule
Risk Factor & Mitigation Strategy \\
\midrule
Cold-Start Anomalies & Initialize with category-level population priors; decay weights exponentially toward personalized model. \\
Data Overfitting & Cross-validate against recent task splits; regularize small user datasets toward global priors. \\
Manual Overrides & Log overrides as explicit feedback features; display model confidence intervals. \\
Data Privacy & Exclude raw task payload content; store only anonymized timestamps, categories, and metrics. \\
\bottomrule
\end{tabularx}
\end{table}

\section{Implementation Timeline}
\begin{table}[htbp]
\centering
\small
\begin{tabularx}{\textwidth}{c X}
\toprule
\textbf{Timeline} & \textbf{Milestone Focus} \\
\midrule
Weeks 1 to 2 & Requirements definition, schema specification, authentication and API core setup. \\
Weeks 3 to 4 & Baseline bias regression implementation and task CRUD operational workflows. \\
Weeks 5 to 7 & Circadian productivity engine, context-switch metrics, and inference pipeline. \\
Weeks 8 to 9 & Deadline proximity integration, auto-rebuilder development, and notification engine. \\
Weeks 10 to 11 & Analytics dashboard implementation, performance tuning, and KPI benchmarking. \\
Week 12 & Comprehensive integration testing, system hardening, and project defense. \\
\bottomrule
\end{tabularx}
\end{table}

\section{Summary}
\textit{TimeLens AI} converts task management into an adaptive, data-driven forecasting ecosystem. By quantifying individual operational bias, circadian performance curves, context-switch overhead, and deadline behavior, the platform replaces manual scheduling with dynamic, statistically backed temporal planning evaluated via clear, empirical KPIs.

\end{document}